\section{总结与延伸}

\subsection{一张表回收整套 Recipe}

下面的表将四个主要系统按输入、输出和失败模式放在一起。它比按论文年份记忆更有用，因为后续任何机器人 foundation model 都必须覆盖这些功能，只是实现可能被合并进同一模型。

\begin{longtable}{@{}L{0.16\textwidth}L{0.22\textwidth}L{0.23\textwidth}L{0.27\textwidth}@{}}
\toprule
系统 & 核心输入 & 核心输出 & 主要失败模式 \\
\midrule
RT-1 & 六帧视觉、语言任务、离线示范 & 离散化实时动作 & 新场景、多因素变化、长时 context 与 dataset coverage \\
SayCan & 高层目标、skill list、affordance scores & 下一项 grounded skill & skill library 太小、affordance 失准、执行失败 \\
Inner Monologue & scene、success、用户反馈、历史计划 & 更新 belief 与下一步 plan & detector / caption 错误、反馈 ontology 不全 \\
DIAL & 小量人类语言、大量 offline trajectories、VLM & 扩展 language labels 与 policy & label noise、行为覆盖不足、组合泛化证据有限 \\
\bottomrule
\end{longtable}

\subsection{最终结论}

第一，机器人 scaling 的核心资源不是单一 episode count，而是任务、环境、语言和 embodiment 的多样覆盖。第二，Transformer 的真正价值之一是 token interface 与高容量 sequence model，但实时预算迫使视觉压缩、短历史和小模型部署。第三，LLM 最适合通过受 grounding 约束的 planner 接入机器人，而不是直接把语言概率当动作。第四，feedback 和 VLM relabeling 能把外部 foundation-model prior 注入系统，却不能凭空创造未采集物理技能。

\begin{importantbox}{最值得带走的一句话}
2023 年的 robotics foundation-model recipe 不是“把机器人接上 LLM”，而是：用大型多样离线数据训练可实时执行的 skill model，用语言连接计划与反馈，用 affordance 保持 grounding，再让外部 foundation models 持续改善语义接口。
\end{importantbox}

\subsection{拓展阅读}

以下资料均在课堂日期前公开，并直接构成本讲证据链。阅读顺序建议是 RT-1 的 skill/data 层、SayCan 的 planner、Inner Monologue 的 feedback，最后用 DIAL 回到数据接口。

\begin{itemize}
  \item \href{https://arxiv.org/abs/2212.06817}{Brohan et al., RT-1: Robotics Transformer for Real-World Control at Scale}：多任务模仿、架构、data scaling 与 ablation。
  \item \href{https://arxiv.org/abs/2204.01691}{Ahn et al., Do As I Can, Not As I Say}：LLM usefulness 与 affordance grounding。
  \item \href{https://arxiv.org/abs/2207.05608}{Huang et al., Inner Monologue}：scene、success 与 active feedback 的闭环规划。
  \item \href{https://arxiv.org/abs/2202.02005}{Jang et al., BC-Z}：语言条件多任务模仿和 zero-shot task generalization。
  \item \href{https://arxiv.org/abs/2211.11736}{Xiao et al., Robotic Skill Acquisition via Instruction Augmentation with Vision-Language Models}：DIAL 与 VLM relabeling。
\end{itemize}

\subsection{课后自检}

\begin{importantbox}{十个自检问题}
\begin{enumerate}
  \item 为什么本讲只能说“towards a robotics foundation model”，不能说已经证明 emergence？
  \item Offline robot learning 如何解耦 data generation 与 consumption，又引入什么 distribution-shift 风险？
  \item RT-1 的 TokenLearner、六帧历史、256 bins 分别解决什么资源约束？
  \item 为什么 task diversity 的减少比同等 episode reduction 更伤 generalization？
  \item SayCan 中 LM score 与 affordance score 各自不能替代什么？
  \item Planning success 与 execution success 为什么必须分开报告？
  \item Inner Monologue 的 passive、active、success feedback 有何不同？
  \item DIAL 为什么能扩展语言覆盖，却不能创造未采集技能？
  \item “少量 label noise okay”依赖哪些 dataset 条件，不能被怎样误读？
  \item 六帧已接近 context limit 时，分钟级任务需要怎样的 memory architecture？
\end{enumerate}
\end{importantbox}

\end{document}
